% Generated draft: requires mathematical and source review.
\documentclass[11pt]{article}
\usepackage[margin=25mm]{geometry}
\usepackage{fontspec}
\setmainfont{Latin Modern Roman}
\usepackage{amsmath,amssymb,mathtools,unicode-math}
\setmathfont{Latin Modern Math}
\usepackage{enumitem,needspace,xurl,hyperref}
\setlength{\parindent}{0pt}
\setlength{\parskip}{5pt}
\setlength{\emergencystretch}{3em}
\newcommand{\N}{\mathbb{N}}
\newcommand{\Z}{\mathbb{Z}}
\newcommand{\Q}{\mathbb{Q}}
\newcommand{\R}{\mathbb{R}}
\newcommand{\C}{\mathbb{C}}
\newcommand{\F}{\mathbb{F}}
\newcommand{\A}{\mathbb{A}}
\newcommand{\PP}{\mathbb{P}}
\newcommand{\E}{\mathbb{E}}
\newcommand{\eps}{\varepsilon}
\newcommand{\dd}{\,\mathrm{d}}
\begin{document}
\section*{UnsolvedMath ID: 30003592}
\textbf{Source problem code:} catalogue entry
\subsection{Definitions and mathematical statement}
Let \(X\) be a smooth toric variety, and write \(n=\dim X\). The source does not specify the ground field or assume that \(X\) is complete or projective. These omissions require clarification before fixing the numerical cycle space and the meaning of movability.

In the usual proper, smooth setting over a fixed field, let \(N_d(X)_{\R}\) denote the real vector space of algebraic \(d\)-cycles modulo numerical equivalence: a cycle is numerically trivial if its intersection with every cycle of complementary dimension has degree zero. Its rational structure is given by classes of cycles with rational coefficients. A customary interpretation of a movable integral \(d\)-dimensional subvariety is one occurring in an algebraic family of subvarieties whose union is Zariski dense in \(X\). The source does not specify the precise family convention; this interpretation requires confirmation.

With the intended numerical equivalence and movability conventions fixed, define
\[
\operatorname{Mov}_d(X)
=
\overline{\operatorname{cone}\bigl\{[V]: V\subseteq X\text{ is a movable integral subvariety of dimension }d\bigr\}}
\subseteq N_d(X)_{\R},
\]
where \(\operatorname{cone}\) consists of finite linear combinations with nonnegative real coefficients, and the closure is taken in the usual topology of the numerical cycle space. A cone is rational polyhedral if it is generated by finitely many rational numerical classes; equivalently, it has the form
\[
\sum_{i=1}^{r}\R_{\geq 0}\alpha_i,
\qquad \alpha_i\in N_d(X)_{\Q}.
\]

\textbf{Question as supplied.} For every smooth toric variety \(X\), is \(\operatorname{Mov}_d(X)\) rational polyhedral for every relevant dimension \(d\)? The natural dimension range is \(0\leq d\leq n\), including its endpoints, but the source states no explicit range or exclusions. No completeness, projectivity, or ground-field hypothesis has been added to the question. The definitions above describe a customary proper-setting interpretation, rather than resolving these missing hypotheses.
\subsection{Short English statement}
The question asks whether the closed cone of numerical classes of movable \(d\)-dimensional subvarieties of a smooth toric variety always has finitely many rational generators. The ground field, completeness or projectivity assumptions, dimension range, and precise movability convention need clarification. The supplied registry labels the problem open; that status has not been independently verified.
\subsection{Sources}
\begin{itemize}
\item Supplied local UnsolvedMath registry record, permanent ID \(30003592\), title: \emph{Polyhedrality of Movable Cycle Cones on Toric Varieties}; exact source problem code: \texttt{catalogue entry}. The collection is \emph{Oberwolfach Reports Open Problems}, extracted from workshop discussions. No individual proposer is supplied.
\item Supplied source citation: \emph{Mini-Workshop: Positivity in Higher-dimensional Geometry: Higher-codimensional Cycles and Newton--Okounkov Bodies} (2018), \emph{Oberwolfach Reports}, pp.\ 2631--2657, DOI \texttt{10.4171/owr/2017/43}. Supplied source URL: \url{https://doi.org/10.4171/owr/2017/43}. The record separately gives proposed year 2017 and the tag \texttt{owr-year:2017}; their relationship to the citation year 2018 is not explained.
\item Supplied literature source, labelled ``original OWR report'': \url{https://doi.org/10.4171/OWR/2017/43}.
\item Supplied literature source, labelled ``Brambilla et al., cones for Mori dream spaces'': \url{https://pure.mpg.de/pubman/item/item_3592150_4/component/file_3644841/Brambilla-Dumitrescu-Postinghel-Sanchez_Duality%20and%20polyhedrality%20of%20cones%20for%20Mori%20dream%20spaces_2025.pdf}.
\item Literature assessment dated \texttt{2026-08-21}: this is supplied dated curation, not independently verified. It labels the problem open and reports polyhedrality in the toric divisor and curve cases, together with later rational polyhedrality results for related ample-in-codimension cones on Mori dream spaces, including toric varieties. The supplied assessment explicitly says these results do not establish the stated rational polyhedrality for movable higher-dimensional subvariety classes in every codimension. The background describes this as dated literature triage derived from the supplied OWR curation and requests independent verification of time-sensitive status and cited claims.
\item The registry supplies the statement-status label \texttt{corrected\_verified}, explained as preservation of mathematical content while notation, typography, and sentence structure were repaired for publication. This supplied editorial label does not establish independent mathematical verification of this draft.
\end{itemize}
\end{document}
